\documentclass[12pt,letterpaper]{article}

\usepackage[utf8]{inputenc}
\usepackage[T1]{fontenc}
\usepackage[spanish,es-nodecimaldot]{babel}
\usepackage{lmodern}
\usepackage{microtype}
\usepackage{geometry}
\usepackage{setspace}
\usepackage{parskip}
\usepackage{booktabs}
\usepackage{longtable}
\usepackage{tabularx}
\usepackage{array}
\usepackage{enumitem}
\usepackage[table]{xcolor}
\usepackage{titlesec}
\usepackage{fancyhdr}
\usepackage{hyperref}
\usepackage{float}

\geometry{top=2.54cm,bottom=2.54cm,left=2.54cm,right=2.54cm}
\doublespacing
\setlength{\parindent}{1.27cm}
\setlength{\parskip}{0pt}
\setlength{\emergencystretch}{3em}

\definecolor{cacao}{HTML}{4B2E20}
\definecolor{selva}{HTML}{2E5D50}
\definecolor{hoja}{HTML}{EAF1ED}
\definecolor{arena}{HTML}{F4F0E9}
\definecolor{gristexto}{HTML}{3D3D3D}

\hypersetup{
  colorlinks=true,
  linkcolor=selva,
  urlcolor=selva,
  citecolor=cacao,
  pdftitle={Análisis técnico de la Sesión 2 - CacaoSelva},
  pdfauthor={Juan Manuel Rengifo Fretel}
}

\pagestyle{fancy}
\fancyhf{}
\fancyhead[L]{\small\textcolor{cacao}{CacaoSelva}}
\fancyhead[R]{\small\textcolor{gristexto}{Construcción de Software II}}
\fancyfoot[C]{\small\textcolor{gristexto}{\thepage}}
\renewcommand{\headrulewidth}{0.4pt}
\renewcommand{\headrule}{\hbox to\headwidth{\color{cacao}\leaders\hrule height \headrulewidth\hfill}}
\setlength{\headheight}{24.02pt}
\renewcommand{\arraystretch}{1.14}

\titleformat{\section}
  {\normalfont\Large\bfseries\color{cacao}}
  {\thesection}{0.75em}{}
\titleformat{\subsection}
  {\normalfont\large\bfseries\color{selva}}
  {\thesubsection}{0.75em}{}
\titlespacing*{\section}{0pt}{2.2ex plus .5ex minus .2ex}{1.1ex}
\titlespacing*{\subsection}{0pt}{1.6ex plus .4ex minus .2ex}{.7ex}

\newcolumntype{Y}{>{\raggedright\arraybackslash}X}
\newcolumntype{P}[1]{>{\raggedright\arraybackslash}p{#1}}
\newcommand{\tablaNota}[1]{\par\smallskip\noindent\textit{Nota.}~#1}
\newcommand{\codigo}[1]{\textbf{#1}}
\newenvironment{referencias}{%
  \list{}{%
    \setlength{\leftmargin}{1.27cm}%
    \setlength{\itemindent}{-1.27cm}%
    \setlength{\itemsep}{0.5em}%
    \setlength{\parsep}{0pt}%
  }%
}{\endlist}

\begin{document}

\begin{titlepage}
\thispagestyle{empty}
\vspace*{1.4cm}
\noindent\textcolor{selva}{\rule{\textwidth}{1.2pt}}
\vspace{0.9cm}
\begin{flushleft}
{\Large\textcolor{selva}{CONSTRUCCIÓN DE SOFTWARE II}}\\[1.25cm]
 {\fontsize{26}{32}\selectfont\bfseries\textcolor{cacao}{Análisis técnico de la Sesión 2}}\\[0.15cm]
 {\fontsize{26}{32}\selectfont\bfseries\textcolor{cacao}{Prueba de concepto CacaoSelva}}
\end{flushleft}
\vspace{0.8cm}
\noindent\textcolor{cacao}{\rule{0.34\textwidth}{4pt}}
\vfill
\begin{flushleft}
\begin{tabular}{@{}ll@{}}
\textbf{Estudiante:} & Juan Manuel Rengifo Fretel\\[0.35cm]
\textbf{Docente:} & Rios Rivera, Carlos Abraham\\[0.35cm]
\textbf{Semestre:} & 2026-II\\[0.35cm]
\textbf{Fecha:} & 26 de septiembre de 2026
\end{tabular}
\end{flushleft}
\vfill
\begin{flushleft}
\small\textcolor{gristexto}{Informe académico de definiciones, alcance y análisis técnico de los requerimientos de la sesión.}
\end{flushleft}
\noindent\textcolor{selva}{\rule{\textwidth}{1.2pt}}
\end{titlepage}

\pagenumbering{arabic}

\section*{Resumen}
\addcontentsline{toc}{section}{Resumen}
Este informe analiza la Sesión 2, titulada \textit{Prueba de concepto en Java con Arquitectura Limpia}, aplicada al caso CacaoSelva. El documento organiza el contenido en tres perspectivas: requerimientos funcionales, requisitos técnicos y arquitectura. El análisis identifica qué comportamiento debe ofrecer la solución, qué tecnologías y restricciones condicionan su implementación y cómo deben estructurarse las dependencias para proteger el dominio. La propuesta se orienta a una prueba de concepto académica compuesta por una API REST, una aplicación Desktop, un Monitor automático y, como actividad complementaria, un cliente web. La conclusión principal es que el éxito del ejercicio no depende únicamente de obtener respuestas correctas, sino de demostrar separación de responsabilidades, inversión de dependencias, manejo controlado de fallas y capacidad de evolución.

\textbf{Palabras clave:} CacaoSelva, Arquitectura Limpia, Java, requisitos, SOLID, API REST.

\clearpage
\tableofcontents
\newpage

\section{Introducción}
La sesión propone construir una prueba de concepto, no un sistema productivo. Esta diferencia es importante: el alcance se concentra en demostrar que varias aplicaciones pueden consultar la misma información y que las reglas del negocio permanecen aisladas de las tecnologías externas. En consecuencia, el proyecto debe privilegiar claridad, trazabilidad y facilidad de explicación sobre características de producción como seguridad avanzada, alta disponibilidad o persistencia distribuida.

El caso utiliza tres lotes ficticios: Ana, con 120,5 kg y estado pendiente; Luis, con 80 kg y estado liquidado; y Rosa, con 95,25 kg y estado pendiente. A partir de ellos se obtiene un total de tres lotes, dos pendientes, uno liquidado y 295,75 kg. Estos datos constituyen la referencia funcional para validar la API, el Desktop y el Monitor.

La arquitectura descrita sigue el principio de que las capas externas pueden conocer a las internas, pero las internas no deben depender de las externas. Este criterio se relaciona con la independencia de frameworks y con la regla de dependencia de la Arquitectura Limpia (Martin, 2017).

\section{Alcance y criterios de análisis}
El análisis se organiza en tres tablas. La primera reúne los requerimientos funcionales y describe el comportamiento observable del sistema. La segunda agrupa los requisitos técnicos y explica las decisiones tecnológicas necesarias para materializar ese comportamiento. La tercera sintetiza la arquitectura propuesta, sus capas, contratos y responsabilidades.

Las tablas no representan código ni una implementación concreta. Funcionan como una especificación técnica para orientar el diseño, la construcción y la validación posterior mediante las pruebas P01 a P08. La distinción entre requisito funcional y técnico permite evitar que una tecnología se confunda con el resultado que el usuario necesita.

\begin{center}
\fcolorbox{selva}{hoja}{%
\begin{minipage}{0.88\textwidth}
\small
\textcolor{cacao}{\textbf{Ruta de lectura}}\\[0.25em]
\textcolor{gristexto}{\textbf{Qué debe hacer} \quad $\longrightarrow$ \quad \textbf{Con qué condiciones} \quad $\longrightarrow$ \quad \textbf{Cómo se organiza}}
\end{minipage}}
\end{center}

\section{Tabla 1. Requerimientos funcionales}
\begin{table}[!ht]
\centering
\caption{Requerimientos funcionales y análisis técnico}
\label{tab:funcionales}
\small\singlespacing
\rowcolors{2}{arena}{white}
\begin{tabularx}{\textwidth}{@{}P{1.25cm}P{4.05cm}Y@{}}
\toprule
\rowcolor{cacao}\color{white}\textbf{Código} & \color{white}\textbf{Requerimiento} & \color{white}\textbf{Análisis técnico}\\
\midrule
RF01 & Obtener todos los lotes mediante \texttt{GET /lotes}. & La API debe retornar los tres registros disponibles en memoria.\\
RF02 & Consultar un lote por identificador. & El sistema debe localizar un lote específico usando su ID.\\
RF03 & Informar cuando un lote no existe. & Un ID válido pero inexistente debe producir HTTP 404 Not Found.\\
RF04 & Rechazar identificadores con formato inválido. & Valores como \texttt{abc} no pueden convertirse en identificadores numéricos y deben producir HTTP 400 Bad Request.\\
RF05 & Rechazar identificadores no permitidos. & Los valores menores o iguales a cero deben considerarse inválidos y producir HTTP 400.\\
RF06 & Permitir consultar lotes desde Desktop. & La aplicación JavaFX debe disponer de una acción visible para iniciar la consulta.\\
RF07 & Mostrar los lotes en una tabla. & La interfaz debe presentar ID, socio, peso y estado de cada lote.\\
RF08 & Informar que una consulta está en proceso. & Mientras se realiza la petición, la interfaz debe mostrar un estado como “Consultando...”.\\
RF09 & Informar cuando la API no está disponible. & Un error de comunicación debe mostrarse claramente y no confundirse con una lista vacía.\\
RF10 & Consultar lotes automáticamente desde el Monitor. & El Monitor debe funcionar sin intervención del usuario.\\
RF11 & Repetir la consulta cada 10 segundos. & Debe utilizarse una planificación periódica controlada.\\
RF12 & Contar lotes pendientes. & El Monitor debe identificar los lotes en estado PENDIENTE y obtener el valor 2.\\
RF13 & Registrar fallas de disponibilidad. & Una interrupción de la API debe registrarse sin provocar el cierre inesperado del Monitor.\\
RF14 & Recuperarse después de una falla temporal. & El Monitor debe continuar ejecutándose y volver a consultar cuando la API se restablezca.\\
RF15 & Detener el Monitor de manera controlada. & El proceso debe permitir finalizar sus recursos y tareas periódicas correctamente.\\
RF16 & Consumir la API desde un cliente web. & El cliente web debe consultar el mismo endpoint \texttt{/lotes}.\\
RF17 & Diferenciar error de conexión y lista vacía. & La aplicación debe distinguir entre una respuesta válida sin registros y una falla de comunicación.\\
\bottomrule
\end{tabularx}
\tablaNota{Los códigos RF corresponden a la especificación funcional de la Sesión 2. La aceptación debe comprobarse mediante respuestas HTTP, observación de la interfaz y registros del Monitor.}
\end{table}

\subsection{Lectura técnica de los requerimientos funcionales}
Los requerimientos RF01 a RF05 establecen el contrato mínimo de la API. Es especialmente relevante diferenciar los errores 400 y 404: el primero representa una solicitud inválida, mientras que el segundo representa una solicitud válida sobre un recurso inexistente. Esta distinción evita respuestas ambiguas y permite que los clientes reaccionen de manera correcta.

Los requerimientos RF06 a RF09 trasladan el contrato al Desktop. La interfaz no debe limitarse a presentar datos: también debe hacer visible el estado de la operación y comunicar la indisponibilidad de la API. El criterio más importante es que una falla de conexión no se presente como “cero lotes”, porque ambos estados tienen significados operativos diferentes.

Los requerimientos RF10 a RF15 definen el comportamiento del Monitor como observador periódico. Cada mensaje “Pendientes: 2” representa una observación del estado actual, no una acumulación de lotes. La recuperación ante fallas exige que el ciclo periódico sobreviva a una excepción temporal y vuelva a intentarlo.

\section{Tabla 2. Requisitos técnicos}
\begin{table}[!ht]
\centering
\caption{Requisitos técnicos y análisis de implementación}
\label{tab:tecnicos}
\small\singlespacing
\rowcolors{2}{arena}{white}
\begin{tabularx}{\textwidth}{@{}P{1.25cm}P{4.05cm}Y@{}}
\toprule
\rowcolor{cacao}\color{white}\textbf{Código} & \color{white}\textbf{Requisito técnico} & \color{white}\textbf{Análisis técnico}\\
\midrule
RT01 & Utilizar Java. & Java proporciona el lenguaje base y las bibliotecas para los componentes.\\
RT02 & Administrar la construcción con Maven. & Maven organiza módulos, dependencias, compilación y pruebas.\\
RT03 & Implementar la API con Spring Boot. & Spring Boot proporciona el servidor HTTP, controladores REST y configuración de la API.\\
RT04 & Implementar Desktop con JavaFX. & JavaFX permite construir la interfaz gráfica y sus controles visuales.\\
RT05 & Implementar el Monitor como proceso Java independiente. & El Monitor puede iniciarse y detenerse separado de la API y Desktop.\\
RT06 & Utilizar HTTP para la comunicación. & Los clientes se comunican con la API mediante solicitudes HTTP.\\
RT07 & Utilizar JSON para el intercambio. & JSON transporta los lotes entre la API y los clientes.\\
RT08 & Ejecutar la API en \texttt{http://localhost:5080}. & La dirección y el puerto deben estar centralizados en la configuración.\\
RT09 & Utilizar \texttt{java.net.http.HttpClient}. & El adaptador HTTP encapsula el consumo mediante el cliente estándar de Java.\\
RT10 & Ejecutar operaciones bloqueantes de forma asíncrona. & Las llamadas HTTP del Desktop no deben ejecutarse en el hilo de JavaFX.\\
RT11 & Consultar aproximadamente cada 10 segundos. & La periodicidad debe implementarse con un planificador controlado.\\
RT12 & Centralizar URL y tiempos de espera. & La configuración no debe repetirse en varios componentes.\\
RT13 & Conectar tecnologías externas mediante adaptadores. & HTTP, JSON, JavaFX y almacenamiento deben mantenerse fuera del núcleo de negocio.\\
RT14 & Utilizar Postman para comprobar la API. & Postman permite validar endpoints, respuestas y errores HTTP.\\
RT15 & Registrar método, URL, código y cuerpo. & Las pruebas deben generar evidencia verificable del comportamiento de la API.\\
\bottomrule
\end{tabularx}
\tablaNota{Los requisitos RT describen restricciones tecnológicas y de operación. No sustituyen los criterios funcionales: indican con qué condiciones debe alcanzarse el comportamiento esperado.}
\end{table}

\subsection{Lectura técnica de los requisitos técnicos}
La combinación de Java y Maven permite organizar el proyecto como una solución multimódulo. Spring Boot y JavaFX cumplen funciones externas diferentes: la primera atiende solicitudes HTTP y la segunda presenta información al usuario. El Monitor, por su parte, se mantiene como un proceso Java independiente para demostrar que varios consumidores pueden usar el mismo contrato.

HTTP y JSON constituyen el límite de comunicación entre aplicaciones. Por ello, los detalles de construcción de solicitudes, conversión de respuestas y tratamiento de códigos HTTP deben concentrarse en un adaptador. La interfaz gráfica y el Monitor solo deberían conocer una operación conceptual de consulta, no los detalles internos de la red.

El requisito RT10 es crítico para la experiencia de usuario. Una operación bloqueante ejecutada directamente en el hilo de JavaFX puede congelar la ventana. La solución debe realizar la consulta fuera de ese hilo y devolver el resultado a la interfaz de forma segura. Del mismo modo, RT11 exige planificación periódica en lugar de un ciclo infinito desorganizado.

\section{Tabla 3. Arquitectura propuesta}
\begin{table}[!ht]
\centering
\caption{Elementos de Arquitectura Limpia y responsabilidades}
\label{tab:arquitectura}
\small\singlespacing
\rowcolors{2}{arena}{white}
\begin{tabularx}{\textwidth}{@{}P{3.05cm}P{4.0cm}Y@{}}
\toprule
\rowcolor{cacao}\color{white}\textbf{Elemento} & \color{white}\textbf{Definición} & \color{white}\textbf{Análisis técnico}\\
\midrule
Arquitectura Limpia & Modelo que protege las reglas de negocio frente a detalles tecnológicos. & Las dependencias deben dirigirse hacia las capas internas.\\
Dominio & Núcleo con las entidades y conceptos principales del negocio. & Debe ser independiente de Spring, JavaFX, HTTP, JSON y bases de datos.\\
Aplicación & Capa que contiene los casos de uso. & Define qué debe hacer el sistema sin conocer cómo se transportan o almacenan los datos.\\
Puertos & Interfaces que representan necesidades de la aplicación. & Permiten depender de abstracciones en lugar de implementaciones concretas.\\
Adaptadores & Componentes que conectan el núcleo con tecnologías externas. & Traducen HTTP, JSON, almacenamiento o interfaz gráfica sin contaminar el dominio.\\
API REST & Adaptador de entrada para solicitudes HTTP. & Recibe peticiones, invoca casos de uso y transforma resultados en respuestas HTTP.\\
Desktop & Interfaz gráfica para la interacción del usuario. & Debe delegar la consulta y no contener reglas de negocio ni detalles HTTP.\\
Monitor & Proceso automático que consulta y registra el estado. & Debe separar planificación, consulta, conteo y registro de errores.\\
Repositorio en memoria & Implementación temporal del acceso a lotes. & Permite demostrar la solución sin una base de datos.\\
Inversión de dependencias & Las capas internas dependen de abstracciones. & Permite reemplazar el almacenamiento con bajo impacto sobre los casos de uso.\\
Separación de responsabilidades & Cada componente posee una función claramente delimitada. & Controladores, casos de uso, repositorios, adaptadores y planificadores evolucionan de forma independiente.\\
Independencia del dominio & El dominio no conoce frameworks ni infraestructura. & Facilita pruebas, mantenimiento y sustitución de tecnologías.\\
Reutilización del núcleo & API, Desktop y Monitor comparten contratos y reglas. & Evita duplicar validaciones y lógica de negocio en cada aplicación.\\
Manejo de errores & Los errores se transforman en respuestas o mensajes controlados. & Se diferencian errores de entrada, recursos inexistentes y fallas de conexión.\\
Recuperación del Monitor & El proceso continúa después de una falla temporal. & Una interrupción de la API no debe detener permanentemente el Monitor.\\
Evidencia de validación & Las pruebas demuestran funcionamiento y arquitectura. & Deben comprobarse P01–P08, incluyendo fallas y detención controlada.\\
\bottomrule
\end{tabularx}
\tablaNota{La dirección conceptual de dependencias es: interfaces externas, adaptadores, aplicación y dominio. Las capas externas pueden conocer a las internas; el dominio no debe conocer a las externas.}
\end{table}

\subsection{Lectura técnica de la arquitectura}
La arquitectura propuesta utiliza el dominio como centro estable. En él se ubican los conceptos que representan el negocio, mientras que la aplicación contiene los casos de uso que coordinan las operaciones. Los puertos expresan las necesidades de esas capas sin imponer una tecnología concreta.

La infraestructura implementa los contratos definidos por las capas internas. En esta prueba de concepto, el repositorio en memoria satisface el acceso a los lotes y el adaptador HTTP satisface la consulta realizada por los clientes. Si posteriormente se incorpora una base de datos, el cambio debe concentrarse en una nueva implementación del puerto, sin modificar las reglas de negocio.

La API, el Desktop y el Monitor son entradas diferentes al mismo problema. La API expone el contrato HTTP; el Desktop convierte el resultado en una experiencia visual; y el Monitor ejecuta observaciones periódicas. Esta separación aplica el principio de responsabilidad única y el principio de inversión de dependencias, dos prácticas centrales para reducir acoplamiento (Martin, 2008).

\section{Trazabilidad y estrategia de validación}
La trazabilidad relaciona cada grupo de requisitos con una forma de comprobación. RF01 a RF05 se validan mediante Postman y sus códigos HTTP. RF06 a RF09 se observan en la interfaz Desktop, incluyendo el estado de consulta y la respuesta ante la caída de la API. RF10 a RF15 se verifican mediante los registros del Monitor, su continuidad tras una falla y su detención voluntaria. RT01 a RT15 se inspeccionan en la configuración, los módulos, los adaptadores y la evidencia de las pruebas.

La validación completa debe incluir los escenarios P01 a P08 descritos en la sesión: listado, consulta individual, lote inexistente, identificadores inválidos, consulta desde Desktop, ejecución del Monitor, interrupción y recuperación de la API y detención controlada del Monitor. La evidencia recomendada incluye capturas de Postman, respuestas JSON, códigos HTTP, mensajes de la interfaz y registros de consola.

\section{Conclusiones}
La Sesión 2 define una prueba de concepto cuyo valor principal es arquitectónico. El sistema debe demostrar que una API, un Desktop y un Monitor pueden consultar la misma información sin duplicar reglas ni acoplar el núcleo a tecnologías externas.

Los requerimientos funcionales establecen el comportamiento observable; los requisitos técnicos fijan las herramientas y restricciones; y la arquitectura organiza las dependencias para proteger el dominio. Estas tres perspectivas son complementarias y deben revisarse conjuntamente.

La solución propuesta es apropiada para un contexto académico porque mantiene el alcance controlado: utiliza datos en memoria, operaciones de solo lectura y ejecución local. No pretende demostrar seguridad, persistencia real, rendimiento, alta disponibilidad ni preparación para producción. Su resultado esperado es una base comprensible y extensible que permita incorporar nuevas implementaciones sin alterar las reglas de negocio.

\section*{Referencias}
\addcontentsline{toc}{section}{Referencias}
\begin{referencias}
\item Martin, R. C. (2008). \textit{Clean code: A handbook of agile software craftsmanship}. Prentice Hall.

\item Martin, R. C. (2017). \textit{Clean architecture: A craftsman's guide to software structure and design}. Prentice Hall.

\item Oracle. (s. f.). \textit{Java documentation}. Recuperado el 26 de septiembre de 2026, de \url{https://docs.oracle.com/en/java/}

\item Spring. (s. f.). \textit{Spring Boot reference documentation}. Recuperado el 26 de septiembre de 2026, de \url{https://docs.spring.io/spring-boot/index.html}

\item OpenJFX. (s. f.). \textit{OpenJFX documentation}. Recuperado el 26 de septiembre de 2026, de \url{https://openjfx.io/}
\end{referencias}

\end{document}
